\subsection{Architecture}
\label{app:software-architecture}

\begin{figure}[htbp]
    \centering
    \begin{tikzpicture}[
        unit/.style={draw, thick, rounded corners=2pt, fill=gray!7, minimum width=1.55cm, minimum height=0.55cm, inner sep=2pt, font=\scriptsize\ttfamily},
        flow/.style={-{Stealth[length=1.6mm]}, thick},
        link/.style={-{Stealth[length=1.4mm]}, densely dashed},
    ]
        \foreach \name/\x/\y in {basis/0/0, operators/1.95/0, problems/3.9/0.45, carleman/3.9/-0.45, assemble/5.85/0.45, ir/5.85/-0.45, classical/7.8/0.45, circuits/7.8/-0.45, verify/9.75/0.45, resources/9.75/-0.45, qsp/11.7/0, pipeline/13.65/0, plotting/15.6/0}
            \node[unit] (\name) at (\x,\y) {\name};
        \foreach \x [count=\k from 0] in {0,1.95,3.9,5.85,7.8,9.75,11.7,13.65,15.6}
            \node[font=\scriptsize, text=gray] at (\x,-1.05) {\k};
        \foreach \x in {0,1.95,3.9,5.85,7.8,9.75,11.7,13.65}
            \draw[flow] (\x+0.82,0) -- (\x+1.95-0.82,0);
        \draw[flow] (assemble.south) -- (ir.north);
        \node[draw, densely dashed, rounded corners=3pt, inner sep=5pt, fit=(basis)(problems)(carleman)(plotting), label={[font=\scriptsize\ttfamily, anchor=south west]north west:src/pihm}] (pkg) {};
        \node[draw, densely dashed, rounded corners=2pt, font=\scriptsize, align=center, inner sep=3pt] (tools) at (13.6,-1.95) {\texttt{tools/} (\texttt{run\_campaign.py}, \texttt{export\_numbers.py}),\\ the notebooks and this thesis's figure scripts};
        \draw[link] ([xshift=4mm]pipeline.south) -- ([xshift=4mm]pipeline.south |- tools.north);
        \draw[link] ([xshift=4mm]plotting.south) -- ([xshift=4mm]plotting.south |- tools.north);
        \node[font=\scriptsize, anchor=west, align=left] at (-0.85,-1.95) {layer \textcolor{gray}{$0$--$8$};\ an arrow points from a unit to the units that may import it,\\ and \texttt{tests/test\_architecture.py} fails on any import against the arrows.};
    \end{tikzpicture}
    \caption[The layers of the \texttt{pihm} package]{\textbf{The \texttt{pihm} package imports in one direction only.} Each unit of \texttt{src/pihm} by its layer, from the Chebyshev basis (0) to plotting (8), with an arrow from a unit to the units that may import it; an architecture test fails on any import against the arrows. The package has no command line: campaigns are run, their summaries collected and the generated numbers exported by the scripts in \texttt{tools/}, outside the package, beside the notebooks and this thesis's figure scripts.}
    \label{fig:pihm-architecture}
\end{figure}

\subsection{Records and provenance}
\label{app:software-records}

\subsection{Campaigns and computational resources}
\label{app:software-campaigns}

\begin{table}[htbp]
    \centering
    \caption[Resource requests of the Setonix campaigns]{\textbf{Resource requests of the Setonix campaigns.} Every request is one node of the \texttt{work} partition with 16 CPUs. A campaign of many short runs is one job, its runs in sequence, for 2 to 18 hours; a campaign of long or large runs is a Slurm array, one task per run, each group of runs asking for the memory and time below. A run that its own resource estimate shows cannot fit its allocation is recorded as infeasible, with the estimate.}
    \label{tab:resource-classes}
    \small
    \begin{tabular}{@{}l l p{0.5\linewidth}@{}}
        \toprule
        Memory & Time & Runs \\
        \midrule
        64 GB & 12 h & the default: every build, verification, preparation and resource estimate not listed below \\
        32 GB & 4--24 h & the circuit simulations (T1) of the linear problems, systems and lifts \\
        64 GB & 24 h & the longest circuit simulations of the descriptor \\
        128 GB & 6--12 h & the larger Navier--Stokes builds and preparations \\
        230 GB & 12 h & the largest Navier--Stokes builds and resource estimates \\
        \bottomrule
    \end{tabular}
\end{table}

\subsection{Tests}
\label{app:software-tests}

\begin{table}[htbp]
    \centering
    \caption[The test tiers of the \texttt{pihm} package]{\textbf{The test tiers of the \texttt{pihm} package.} What each tier tests and when it runs.}
    \label{tab:test-tiers}
    \small
    \begin{tabular}{@{}l p{0.48\linewidth} p{0.26\linewidth}@{}}
        \toprule
        Tier & Content & Run \\
        \midrule
        Unit and property & operator identities over random $n$ and $x$; primitives; arithmetic; IR; filters; phases & every commit \\
        Golden & fixtures from the frozen legacy code and Wu et al.'s printed $n = 2$ matrices & every commit \\
        Reproduction & Wu et al.'s targets, both variants & every commit \\
        Integration & every problem, both forms, smallest $n$, all stages & every commit \\
        Documentation & Sphinx build, doctests, spelling & every commit \\
        Claims & one test per claim; writes the claims report & at gates \\
        Slow & larger-$n$ verification; circuit-simulation ceilings & at gates, on Setonix \\
        \bottomrule
    \end{tabular}
\end{table}

\subsection{Regenerating the results}
\label{app:software-regenerate}

\subsection{Environment}
\label{app:software-environment}
